\documentclass[twocolumn,10pt]{article}
\usepackage{amsmath}
\usepackage{graphicx}

\title{Two-Column Article Example}
\author{Author}
\date{2026}

\begin{document}
\maketitle

\begin{abstract}
This is a two-column article demonstrating the \texttt{twocolumn} class option combined with full-width abstract and title.
\end{abstract}

\section{Introduction}
The two-column format is standard in academic publishing, providing denser information per page.

\section{Method}
Our method consists of three steps: preprocessing, modelling, and postprocessing. Each step handles distinct subtasks.

\section{Column-Spanning Figure}

\begin{figure*}[t]
\centering
Figure spanning both columns would go here.
\caption{A wide figure that crosses both columns using \texttt{figure*}.}
\label{fig:wide}
\end{figure*}

\section{Results}
Performance data in Figure~\ref{fig:wide}.

\section{Discussion}
We observe that the two-column format affects line-break decisions and table widths.

\section{Conclusion}
Future work extends this to three-column variants.
\end{document}
